\documentclass[UTF8,11pt]{ctexart}

% 本稿用于中文讨论；正式投稿以英文稿 paper/main.tex 为准。
% 外部结果仅在相同软件和候选协议下复现后才可作为本地数值比较；投稿时
% 另行执行 venue 的页数、格式、披露和匿名性检查。

\usepackage[margin=1in]{geometry}
\usepackage{amsmath,amssymb}
\usepackage{array}
\usepackage{booktabs}
\usepackage{graphicx}
\usepackage{multirow}
\usepackage{xcolor}
\usepackage{tikz}
\usetikzlibrary{arrows.meta,calc,fit,positioning}
\usepackage{hyperref}

\hypersetup{
  hidelinks,
  pdfauthor={},
  pdfsubject={},
  pdfkeywords={串联质谱, 分子检索, 学习排序}
}

\newcommand{\method}{SpecEmbedding-Rerank}
\newcommand{\base}{SpecMolAlign}
\newcommand{\zs}{\mathbf{z}_s}
\newcommand{\zm}{\mathbf{z}_m}
\newcommand{\candidates}{\mathcal{C}}

\title{固定候选池上的 MS/MS 分子非生成式监督重排序}
\author{匿名作者}
\date{}

\begin{document}

\maketitle

\begin{abstract}
从串联质谱中检索分子结构具有较大挑战，因为质量或分子式约束后的候选池中仍包含大量
化学上相近的 hard negatives。本文研究监督式、非生成式的第二阶段重排序能否改善固定候选
列表，而不生成新的结构。\method{} 首先对齐峰序列谱图编码器和分子图编码器，并复用至多
256 个所提供候选的保存嵌入；随后先逐候选 coarse 打分，再在 coarse top-40 上计算谱图条件
的相对候选偏好。相对偏好由有序 pair 分数与其反向分数之差构成，因此具有反对称性，并对
候选列表置换保持等变；新方法不使用检索 rank，也不在训练时插入正例。我们将其与参数近似
匹配、无 rank 的 pointwise 对照比较。在 MassSpecGym 的质量和分子式候选池上，保存嵌入的
cache-local 评价显示两种监督重排序器都改善了固定跨模态检索器，而 pointwise 在主要 R@1
和 MRR 汇总上略高。因此，本文将相对分支作为形式化的候选交互机制并研究其经验边界，而
不声称关系模块或外部基线具有普遍优势。本文方法是闭集、非生成式方法：它只能改善已有
候选的顺序，不能恢复候选池之外的分子。
\end{abstract}

\noindent\textbf{关键词：}
生物信息学数据挖掘；串联质谱；分子检索；学习排序；多模态表示学习

\section{引言}

串联质谱（MS/MS）被广泛用于代谢组学、天然产物发现和相关生物医学任务中的小分子
鉴定。实验谱图只是分子结构在特定实验条件下的部分观测：不同分子可能产生相似碎片，
同一分子也可能因仪器、加合物和碰撞能量不同而产生差异明显的谱图。因此，结构注释
天然属于一个大规模数据挖掘问题，即根据查询谱图，从受化学条件约束的候选池中检索并
排序正确分子。

MassSpecGym 提供了结构不重叠的数据划分，以及标准化的质量约束和分子式约束检索
候选池 \cite{bushuiev2024massspecgym}。现有方法主要尝试缩小谱图和分子之间的模态差距。
JESTR 使用 contrastive multiview coding 学习谱图--分子联合空间，并直接按余弦相似度
排序候选 \cite{kalia2025jestr,tian2020cmc}。GLMR 则先检索上下文候选，再利用条件
语言模型生成 refined molecule，最后根据分子--分子相似度重排候选
\cite{zhang2026glmr}。

这两类方法仍留下一个重要的数据挖掘问题。逐候选余弦打分并未在已检索集合上显式训练
第二阶段目标；生成式检索则引入自回归解码器，最终排序依赖生成结构的质量和有效性。因此，
本文回答两个相互关联的问题：监督第二阶段排序器能否区分 hard candidates，以及谱图条件的
候选相对比较是否比容量匹配的 pointwise 打分提供额外收益。前一个问题检验学习式列表重排
本身是否有价值，后一个问题检验该收益是否必须依赖候选交互。在每组配对比较中，候选池和
第一阶段表示保持固定，使比较集中于第二阶段排序机制；跨 alignment 矩阵则单独描述不同
保存的第一阶段 checkpoint 带来的敏感性。

本文基于 SpecEmbedding 提出的峰序列谱图架构 \cite{xiong2025specembedding} 构建
\method{} 两阶段框架。第一阶段从零训练该谱图塔，并将其与 GINE 分子图编码器对齐；第二
阶段先对完整供给池使用绝对分支，再对 coarse top-$K_r$ 子集使用谱图条件的候选关系。
有序 pair 偏好通过反向 pair 分数相减得到，从而无需 rank embedding 即可满足反对称性。训练
只使用自然进入候选池的正例，代码不通过插入目标修复截断列表。具有相同绝对分支的
pointwise 模型是主要容量控制，早期 candidate-set Transformer 仅作为历史参考。

本文贡献如下：

\begin{itemize}
  \item 在 MassSpecGym 候选池上提出直接、闭集、非生成式的第二阶段 learning-to-rank 框架，
  采用完整候选池 coarse 打分和谱图条件的相对候选比较。
  \item 定义无 rank、具有 pair 反对称性、对候选置换保持等变且不执行正例强制插入的相对分支，
  并设置参数近似匹配的 pointwise 对照。
  \item 在质量和分子式候选协议、三个保存的 alignment checkpoint 以及 30 次 relative 模块消融矩阵
  上系统考察该框架：监督重排序改善基础检索器，但当前证据尚不能建立相对候选交互相对
  pointwise 打分的独立优势。
\end{itemize}

\paragraph{AI 使用披露。}
OpenAI Codex 在本稿所有章节中辅助文字起草和语言修改，并在项目过程中辅助代码编辑、
实验编排和结果一致性审计。文中实验观测和数值来自所述软件流水线，而非 AI 模型生成。
作者已审阅全部 AI 辅助文字与代码，核验技术论述、引用、保存的实验工件和报告结果，
并对稿件的准确性、完整性和原创性承担全部责任。

\section{相关工作}

\subsection{结构注释范式}

计算式 MS/MS 注释方法的检索对象和排序证据各不相同，候选分数修正也并非本文首次提出。
MetFusion 融合 in-silico fragmentation 与谱图库证据 \cite{gerlich2013metfusion}；
MS2Query 用五特征随机森林重排 top-2000 MS2DeepScore 匹配
\cite{dejonge2023ms2query}；LC-MS2Struct 将 MS/MS scorer 与保留顺序结合，联合排序
同一 LC 分析中的候选 \cite{bach2022lcms2struct}。本文则基于冻结的谱图--分子表示，对
MassSpecGym 的查询特定候选池排序，不依赖谱图库邻居或保留时间证据
\cite{bushuiev2024massspecgym}。

在结构数据库检索中，CSI:FingerID 和 MIST 从谱图预测分子指纹
\cite{duhrkop2015csifingerid,goldman2023mist}。CMSSP 和 JESTR 学习谱图--分子联合空间；
MVP 进一步在分子图、指纹、单条谱图和 consensus spectra 间进行对比对齐
\cite{chen2024cmssp,kalia2025jestr,chen2026mvp}。这些方法使用预测指纹分数或跨模态
相似度排序候选。SpecEmbedding 研究谱图--谱图表示学习 \cite{xiong2025specembedding}；
本文只复用其峰序列架构，从零开始将其与分子图对齐，并学习已检索 hard-candidate 列表
内部的第二阶段顺序。

其他范式采用不同的中间对象：MassFormer 根据分子图预测串联质谱
\cite{young2024massformer}，MSNovelist 根据 MS/MS 推导的指纹生成结构
\cite{stravs2022msnovelist}，GLMR 则以预检索候选为条件生成分子，再与候选比较
\cite{zhang2026glmr}。本文的闭集方法不解码结构，也不能返回第一阶段列表外的分子；
它直接优化已有候选的分数。

\subsection{学习排序与候选集合建模}

learning-to-rank 使用候选偏好或完整列表定义监督，而不是把相似度排序作为固定后处理。
RankNet 学习 pairwise 偏好，ListNet 则对完整排序列表建立概率模型
\cite{burges2005ranknet,cao2007listnet}。排序函数还可利用竞争候选之间的关系。
Set Transformer 为集合输入提供 attention 模块 \cite{lee2019settransformer}；SetRank
用 self-attention 建模文档间交互 \cite{pang2020setrank}。本文将这些原则用于缓存的
跨模态嵌入，结合显式谱图--候选 pair 特征以及第一阶段分数和排名先验。
表~\ref{tab:method_comparison_cn} 将这一查询特定、非生成式第二阶段与若干相关
MassSpecGym 方法并列定位。本文贡献并非 reranking 或 self-attention 本身；残差打分和
候选交互是接受评估的设计选项，而不是预先假定的增益来源。

\begin{table}[t]
  \centering
  \caption{本文方法与若干相关 MassSpecGym 方法的机制对比。}
  \label{tab:method_comparison_cn}
  \small
  \begin{tabular}{>{\raggedright\arraybackslash}p{1.4cm}
                  >{\raggedright\arraybackslash}p{3.1cm}
                  >{\raggedright\arraybackslash}p{2.8cm}
                  >{\raggedright\arraybackslash}p{1.8cm}
                  >{\raggedright\arraybackslash}p{3.0cm}}
    \toprule
    方法 & 最终排序信号 & 候选间交互 & 生成分子 & 优化阶段 \\
    \midrule
    JESTR & 谱图--候选余弦相似度 & 推理时无 & 否
      & InfoNCE；可选的后期候选正则 \\
    GLMR & 生成分子--候选余弦相似度 & 检索候选作为生成器条件 & 是
      & dual-path InfoNCE 加自回归生成 \\
    本文 & pair 特征和检索先验预测的残差分数
      & 谱图条件的反对称候选对交互 & 否
      & 跨模态 InfoNCE 加 listwise/pairwise 重排序 \\
    \bottomrule
  \end{tabular}
\end{table}

\section{问题定义}

记 $s$ 为一条 MS/MS 谱图，$m^+$ 为真实分子，$\mathcal{D}_s$ 为查询对应的质量约束
或分子式约束的所提供候选池。基础检索器为每个 $m\in\mathcal{D}_s$ 计算

\begin{equation}
  b(s,m)=\operatorname{sim}(f_s(s),f_m(m)),
\end{equation}

并返回 top-$K$ 列表 $\candidates_s=(m_1,\ldots,m_K)$。重排序器学习
$g(s,m_i,\candidates_s)$，使 $m^+$ 的排名尽可能靠前。由于重排序器无法返回
$\candidates_s$ 之外的分子，端到端 Recall@$k$ 的上界为

\begin{equation}
  U_K=\frac{1}{N}\sum_{n=1}^{N}
  \mathbb{I}[m_n^+\in\candidates_{s_n}].
  \label{eq:upper_bound_cn}
\end{equation}

对已进入 top-$K$ 的正例，定义条件排序成功率

\begin{equation}
  A_{k\mid K}=
  \frac{\sum_{n=1}^{N}\mathbb{I}[m_n^+\in\candidates_{s_n},\,
  \operatorname{rank}_g(m_n^+)\le k]}
  {\sum_{n=1}^{N}\mathbb{I}[m_n^+\in\candidates_{s_n}]},
  \qquad
  \mathrm{Recall@}k=U_K A_{k\mid K}\le U_K.
  \label{eq:coverage_decomposition_cn}
\end{equation}

这是指标定义的恒等分解。在固定 top-$K$ cache 时，reranker 只能改变列表内的
$A_{k\mid K}$，不能改变第一阶段覆盖率 $U_K$。因此，我们同时报告 $U_K$ 和面向
全部查询的排序指标；如果正例不在 top-$K$ 中，该查询记为零，而不是只在成功召回的
子集上报告条件准确率。

\section{方法}

\subsection{总体框架}

框架包含跨模态对齐、完整候选池检索和相对重排序三个步骤；后者采用监督的 coarse-to-fine
流程。训练 reranker 时冻结谱图和分子编码器，因此两个模态的嵌入可以离线计算并重复使用。

\begin{figure}[t]
  \centering
  % Shared vector overview for paper/main.tex and paper/main_cn.tex.
% Usage: \methodoverview{0} for English, \methodoverview{1} for Chinese.
\newcommand{\methodoverview}[1]{%
  \begingroup
  \def\molabel##1##2{\ifnum#1=0\relax##1\else##2\fi}%
  \begin{tikzpicture}[
    x=1cm,
    y=1cm,
    font=\scriptsize,
    >=Latex,
    moPanel/.style={
      draw=black!55,
      rounded corners=2pt,
      line width=0.55pt,
      fill=black!2
    },
    moAlignPanel/.style={moPanel, fill=blue!5},
    moRetrievePanel/.style={moPanel, fill=orange!6},
    moRerankPanel/.style={moPanel, fill=green!5},
    moBox/.style={
      draw=black!65,
      rounded corners=1.5pt,
      line width=0.5pt,
      fill=white,
      align=center,
      inner sep=1.7pt,
      minimum height=0.50cm
    },
    moInput/.style={moBox, fill=blue!8},
    moFrozen/.style={moBox, draw=blue!55!black, densely dotted, fill=blue!4},
    moTrain/.style={moBox, draw=orange!75!black, dashed, fill=orange!7},
    moChoice/.style={moBox, draw=green!45!black, fill=green!7},
    moOutput/.style={moBox, draw=purple!55!black, fill=purple!5},
    moArrow/.style={-{Latex[length=1.2mm,width=0.8mm]}, line width=0.65pt},
    moTrainArrow/.style={
      -{Latex[length=1.2mm,width=0.8mm]},
      dashed,
      draw=orange!80!black,
      line width=0.6pt
    },
    moSmall/.style={font=\scriptsize, align=center},
    moTitle/.style={font=\footnotesize\bfseries, align=center}
  ]
    \path[use as bounding box] (0,0) rectangle (12,5.60);

    % Stage panels remain distinguishable in grayscale through borders and titles.
    \node[moAlignPanel, anchor=south west, minimum width=3.95cm,
      minimum height=2.70cm] at (0,2.90) {};
    \node[moRetrievePanel, anchor=south west, minimum width=7.90cm,
      minimum height=2.70cm] at (4.10,2.90) {};
    \node[moRerankPanel, anchor=south west, minimum width=12.00cm,
      minimum height=2.75cm] at (0,0) {};

    \node[moTitle] at (1.975,5.31)
      {\molabel{1. Cross-modal alignment}{1. 跨模态对齐}};
    \node[moTitle] at (8.05,5.31)
      {\molabel{2. Top-$K$ retrieval and caching}{2. Top-$K$ 检索与缓存}};
    \node[moTitle] at (6.00,2.51)
      {\molabel{3. Supervised residual reranking}{3. 监督残差重排序}};

    % Stage 1: spectrum and molecule encoders in a common normalized space.
    \node[moInput, text width=1.20cm] (spectrum) at (0.68,4.47)
      {\molabel{MS/MS\\peaks}{MS/MS\\峰序列}};
    \node[moBox, text width=1.65cm] (specenc) at (2.24,4.47)
      {\molabel{Peak\\Transformer}{谱图\\Transformer}};
    \node[moBox, circle, inner sep=1.3pt, minimum size=0.52cm] (zs) at (3.55,4.47)
      {$\mathbf z_s$};

    \node[moInput, text width=1.20cm] (molecule) at (0.68,3.42)
      {\molabel{Molecule\\graphs}{候选\\分子图}};
    \node[moBox, text width=1.65cm] (molenc) at (2.24,3.42)
      {\molabel{GINE + size\\features}{GINE +\\图规模特征}};
    \node[moBox, circle, inner sep=1.3pt, minimum size=0.52cm] (zm) at (3.55,3.42)
      {$\mathbf z_{m,i}$};

    \draw[moArrow] (spectrum) -- (specenc);
    \draw[moArrow] (specenc) -- (zs);
    \draw[moArrow] (molecule) -- (molenc);
    \draw[moArrow] (molenc) -- (zm);
    \draw[<->, dashed, draw=orange!80!black, line width=0.6pt]
      (zs.south) -- node[left, xshift=-1pt, moSmall, text=orange!70!black,
      fill=blue!5, inner sep=0.5pt]
      {$\mathcal L_{\mathrm{align}}$} (zm.north);

    % Stage 2: score the supplied candidate library and retain a reusable list.
    \node[moBox, text width=1.48cm] (base) at (5.02,4.30)
      {\molabel{Cosine base score}{余弦基础分数}\\$b_i=\mathbf z_s^\top\mathbf z_{m,i}$};
    \node[moBox, text width=1.12cm] (topk) at (6.60,4.30)
      {\molabel{Sort and keep}{排序并保留}\\top-$K=40$};
    \node[moOutput, text width=1.70cm] (cache) at (8.25,4.30)
      {\molabel{Cached list}{缓存候选列表}\\$\{\mathbf z_{m,i},b_i,r_i\}$, $i=1{:}K$};
    \node[moFrozen, text width=2.05cm] (frozen) at (10.55,4.30)
      {\molabel{Encoders frozen;\\embeddings reused}{编码器冻结；\\嵌入重复使用}};
    \node[moSmall, text width=2.30cm] (protocol) at (10.55,3.38)
      {\molabel{MassSpecGym-supplied\\mass/formula candidates}
      {MassSpecGym 提供的\\质量/分子式候选}};

    \draw[moArrow] (zs.east) -| ([yshift=1.4mm]base.west);
    \draw[moArrow] (zm.east) -| ([yshift=-1.4mm]base.west);
    \draw[moArrow] (base) -- (topk);
    \draw[moArrow] (topk) -- (cache);
    \draw[densely dotted, line width=0.55pt, draw=blue!55!black]
      (cache.east) -- (frozen.west);

    % Stage 3: matched pointwise/candidate-set variants and a shared residual path.
    \node[moBox, text width=3.12cm, minimum height=0.78cm] (features) at (1.72,1.55)
      {\molabel{Pair features}{候选配对特征}\\[-1pt]
       $\mathbf x_i=[\mathbf z_s,\mathbf z_{m,i},\mathbf z_s\!\odot\!\mathbf z_{m,i},$\\[-2pt]
       $|\mathbf z_s-\mathbf z_{m,i}|,b_i,\mathbf e(r_i)]$};
    \node[moBox, text width=0.90cm] (mlp) at (3.98,1.55)
      {\molabel{Shared}{共享}\\MLP\\$\mathbf h_i$};
    \node[moChoice, text width=1.70cm] (pointwise) at (5.62,1.86)
      {\textbf{Pointwise}\\\molabel{no list context}{无列表上下文}};
    \node[moChoice, text width=1.90cm] (setaware) at (5.62,1.16)
      {\molabel{\textbf{Candidate-set}\\\textbf{Transformer}}
      {\textbf{候选集合}\\\textbf{Transformer}}\\\molabel{self-attention}{自注意力}};
    \node[moBox, text width=0.90cm] (head) at (7.20,1.55)
      {\molabel{Score head}{打分头}\\$\Delta_i$};
    \node[moBox, text width=1.42cm] (residual) at (8.82,1.55)
      {\molabel{Residual score}{残差分数}\\$q_i=\alpha b_i+\Delta_i$};
    \node[moOutput, text width=1.72cm] (output) at (10.82,1.55)
      {\molabel{Sort by $q_i$}{按 $q_i$ 排序}\\\molabel{same top-$K$ molecules}{仍为同一 top-$K$ 集合}};

    \draw[moArrow] (cache.south) -- (8.25,2.82) -- (0.02,2.82) |- (features.west);
    \draw[moArrow] (features) -- (mlp);
    \draw[moArrow] (mlp.east) -- ++(0.12,0) |- (pointwise.west);
    \draw[moArrow] (mlp.east) -- ++(0.12,0) |- (setaware.west);
    \draw[moArrow] (pointwise.east) -- ++(0.10,0) |- ([yshift=1.2mm]head.west);
    \draw[moArrow] (setaware.east) -- ++(0.10,0) |- ([yshift=-1.2mm]head.west);
    \draw[moArrow] (head) -- (residual);
    \draw[moArrow] (residual) -- (output);

    % Base-score skip and training/test protocol annotations.
    \draw[moArrow, draw=blue!60!black]
      (cache.south east) -- ++(0.12,-0.10) -| node[pos=0.88, right,
      moSmall, text=blue!60!black, fill=green!5, inner sep=1pt]
      {\molabel{base path}{基础路径}: $\alpha b_i$} (residual.north);
    \node[moTrain, text width=4.18cm, minimum height=0.40cm, moSmall] (trainnote) at (5.05,0.45)
      {\molabel{Training only: $\mathcal L_{\mathrm{list}}+\lambda\mathcal L_{\mathrm{pair}}$;
       missing-$m^+$ insertion}
      {仅训练：$\mathcal L_{\mathrm{list}}+\lambda\mathcal L_{\mathrm{pair}}$；
       插入缺失正例 $m^+$}};
    \node[moFrozen, text width=2.95cm, minimum height=0.40cm, moSmall] (testnote) at (10.15,0.45)
      {\molabel{Val./test: reorder only; no insertion/generation}
      {验证/测试：只重排；不插入、不生成}};
    \draw[moTrainArrow] (residual.south west) |- (trainnote.east);
    \draw[densely dotted, line width=0.55pt, draw=blue!55!black]
      (output.south) -- (testnote.north);
  \end{tikzpicture}%
  \endgroup
}

  \methodoverview{1}
  \caption{\method{} 总体流程。面板 1 对齐谱图 Transformer 和分子 GINE；面板 2
  使用冻结、缓存的嵌入构造完整候选池；面板 3 先逐候选 coarse 打分，再在 coarse top-40
  上计算谱图条件的相对候选偏好，历史 Transformer 仅作为历史参考。虚线表示仅训练时使用的
  操作；任何阶段均不执行正例强制替换，整个流程不生成分子。}
  \label{fig:overview_cn}
\end{figure}

\subsection{跨模态基础检索器}

谱图编码器沿用 SpecEmbedding 的峰序列架构 \cite{xiong2025specembedding}，结合正弦
$m/z$ 特征和 Transformer 处理峰的 $m/z$ 与强度序列。在本文实现中，每条谱图最多表示为
100 个 token：第一个 token 是强度设为 2 的前体 $m/z$，其后为按强度选出的至多 99 个
碎片峰。被选碎片恢复其在源数组中的顺序，强度除以所选碎片的最大强度。谱图编码器只接收
这些 $m/z$--强度 token；加合物、仪器和碰撞能量元数据均不是显式输入。本文从零初始化该
谱图塔，并将它用于谱图--分子而非谱图--谱图对齐。分子编码器是基于原子和化学键特征的 GINE
\cite{hu2020pretraininggnn}，并将
mean-pooled 图特征与显式图规模特征拼接。两个投影头将编码结果映射到统一的 $d$ 维空间：

\begin{equation}
  \zs=\frac{p_s(f_s(s))}{\lVert p_s(f_s(s))\rVert_2},
  \qquad
  \zm=\frac{p_m(f_m(m))}{\lVert p_m(f_m(m))\rVert_2}.
\end{equation}

基础模型使用双向跨模态多正例对比损失。具有相同分子标签的谱图和分子均作为正例，
从而避免把同一分子的重复谱图误当成负例。对大小为 $B$ 的 minibatch，记 $P(i)$ 为
与样本 $i$ 具有相同标签的索引集合。使用可学习逆温度尺度 $a$，基础 logits 为

\begin{equation}
  \ell_{ij}=a\,\mathbf{z}_{s_i}^{\top}\mathbf{z}_{m_j}.
\end{equation}

spectrum-to-molecule 项为

\begin{equation}
  \mathcal{L}_{s\rightarrow m}
  =-\frac{1}{B}\sum_{i=1}^{B}\frac{1}{|P(i)|}
  \sum_{j\in P(i)}
  \log\frac{\exp(\ell_{ij})}{\sum_{k=1}^{B}\exp(\ell_{ik})}.
\end{equation}

molecule-to-spectrum 项对 $\ell^\top$ 使用同一表达式，两项取均值得到 alignment loss。
实现令 $a=\exp(u)$ 并将其上限截为 100。推理时，归一化点积
$b_i=\zs^\top\mathbf{z}_{m_i}$ 作为基础分数。

\subsection{候选构造与训练协议}

MassSpecGym 提供按目标分子键控的候选列表；每条 query 复用其目标分子对应的列表。基础
检索器对最多 256 个完整候选打分。新的 coarse-to-fine reranker 先对完整池逐候选打分，
再对 coarse top-$K_r=40$ 候选执行关系精排。256 和 40 都是本修订中预先固定的设计与计算
预算选择；评价脚本可在不插入正例的情况下截断已保存的完整池。

训练只使用正例自然出现在完整 256 候选池中的 query；没有自然正例的 query 保留在覆盖率
统计中，但不进入 reranker 损失。训练时打乱候选顺序，并同步移动候选嵌入和基础分数。新模型
完全不使用 rank embedding，cache 兼容字段 \texttt{base\_ranks} 也不进入其 forward 路径。因此，
任何阶段都不会把正例插入截断列表；较小评价截断中的缺失正例被视为覆盖率 miss。

\subsection{谱图条件的相对候选重排序器}

对候选 $m_i$ 保留绝对匹配特征，但删除检索 rank：

\begin{equation}
  \mathbf{x}_i=[\zs,\ \mathbf{z}_{m_i},\ \zs\odot\mathbf{z}_{m_i},
  |\zs-\mathbf{z}_{m_i}|,\ b_i].
  \label{eq:relative_features_cn}
\end{equation}

绝对分支的 MLP 产生 coarse 分数 $c_i=f_{\rm abs}(\mathbf{x}_i)+\alpha b_i$。
对 coarse top-$K_r$ 候选，令 $u_i=W_m\mathbf{z}_{m_i}$、$v=W_s\zs$，并用谱图条件、
有向候选差异 $u_i-u_j$、其绝对值、基础分数差 $b_i-b_j$ 以及 $u_i,u_j$ 的余弦相似度
构造有序关系向量 $\rho_{ij}$。共享 MLP 定义有向偏好：

\begin{equation}
  p_{ij}=g(\rho_{ij})-g(\rho_{ji}),
  \qquad p_{ij}=-p_{ji}.
  \label{eq:antisymmetry_cn}
\end{equation}

关系权重以谱图为条件，并在有效的非自身候选上归一化。令 $w_{ij}\ge0$，最终分数为

\begin{equation}
  q_i=c_i+\beta\sum_{j\ne i}w_{ij}p_{ij}.
  \label{eq:relative_score_cn}
\end{equation}

关系分支只在 coarse top-$K_r$ 上计算，结果再散射回完整候选池。所有运算在候选维度共享，
并使用附着在候选上的 mask；因此共同置换候选及其特征会以同样方式置换输出分数，但 coarse
top-$K_r$ 选择中的 tie 处理可能影响这一过程。式~\ref{eq:antisymmetry_cn} 给出与参数取值
无关的原始 pair 偏好反对称性；式~\ref{eq:relative_score_cn} 中的加权贡献不要求反对称，
因为 $w_{ij}$ 与 $w_{ji}$ 可以不同。模型不直接使用原始谱峰、
SMILES token、前体质量、分子式或采集元数据，也没有 rank embedding。

\paragraph{结构解释。}
反向 pair 分数之差使一个有序候选对只有一个偏好方向：交换 $i$ 和 $j$ 时，原始 pair 偏好
$p_{ij}$ 只改变符号。这种归纳偏置适合表达“对当前谱图而言候选 $i$ 优于候选 $j$”这样的局部判断，也与训练
目标中的 pairwise 项相容。候选维度共享运算则在列表层面施加相应的置换等变性。但这些是函数
类别的结构性质，不是性能保证：成对反对称本身不能排除三个或更多候选之间的循环偏好，不能
证明 hard-negative discrimination 更好，也不能推出相对模型优于 pointwise 对照。这些问题仍
需要未来经授权的独立实验检验。

当前实现提供显式的 \texttt{pair\_mode} 控制。论文报告的 relative 模型使用 antisymmetric 模式，
即实现式~\ref{eq:antisymmetry_cn}；directed 模式保留同一个有序偏好网络，但省略反向 pair 相减，
作为受控的实现入口。本文没有报告 directed 训练运行，因此该接口不构成新的性能结果。旧的布尔配置
参数仅用于兼容；旧 checkpoint 缺少新字段时仍回退到原有的 antisymmetric 语义。

参数近似匹配的 pointwise 对照使用同一绝对分支，仅移除相对聚合。隐藏宽度 512 时，相对模型有
1,485,189 个可训练参数，无 rank 的 pointwise 对照有 1,478,690 个，差异为 0.44\%。通用
candidate-set Transformer 作为历史参考保留，不再作为本文提出的机制。两种模型都不生成分子。

\subsection{历史 Candidate-Set Transformer 参考}

为提供背景，保留原先含 rank embedding 和候选 self-attention 的 Candidate-Set Transformer
\cite{vaswani2017attention} 作为历史参考。其历史 forcing 协议和 rank-aware 配置不属于本文正式方法；
详细区别见 Appendix~B，且不与当前主表作数值比较。

\subsection{训练目标}

记 $\mathcal{B}_{L}$ 为一个 minibatch 中具有标签的查询集合，$V_n$ 为查询 $n$ 的有效、
非 padding 候选索引集合，$y_n\in V_n$ 为其正例索引。实现首先在所有有标签查询上平均
listwise cross-entropy：

\begin{equation}
  \mathcal{L}_{\mathrm{list}}
  =-\frac{1}{|\mathcal{B}_{L}|}
  \sum_{n\in\mathcal{B}_{L}}
  \log\frac{\exp(q_{ny_n})}{\sum_{i\in V_n}\exp(q_{ni})}.
\end{equation}

此外，每个正例都与同一检索列表中的其他全部有效候选进行比较。实现不是先对每个可变长
列表分别平均，而是在整个 minibatch 的所有有效正负对上统一取均值：

\begin{equation}
  \mathcal{L}_{\mathrm{pair}}
  =\frac{
  \sum_{n\in\mathcal{B}_{L}}\sum_{i\in V_n\setminus\{y_n\}}
  \log\!\left(1+\exp(q_{ni}-q_{ny_n}+\gamma)\right)}
  {\sum_{n\in\mathcal{B}_{L}}\left(|V_n|-1\right)}.
\end{equation}

总损失为

\begin{equation}
  \mathcal{L}
  =\mathcal{L}_{\mathrm{list}}
  +\lambda_{\mathrm{pair}}\mathcal{L}_{\mathrm{pair}}
  +\lambda_{\mathrm{spec}}\mathcal{L}_{\mathrm{spec}}.
\end{equation}
对同一 minibatch 中循环错配的谱图 $s'$，谱图依赖项比较真实候选在正确和错误 query 下的分数：

\begin{equation}
  \mathcal{L}_{\mathrm{spec}}
  =\operatorname{softplus}\!\left(
  q(s',m^+)-q(s,m^+)+\delta\right).
  \label{eq:spectrum_dependency_cn}
\end{equation}

这是 minibatch 内的依赖约束，不是因果识别结论。
其中 listwise 项是一正例、多候选的 top-one likelihood，不规定负例之间的完整全序。
两个损失都对同一查询内全部 $q_i$ 加上任意相同常数保持不变。因此，原始 $q_i$ 是相对
排序分数而非概率；本文既未执行也未评价概率校准。

对嵌入宽度 $d$，完整池 coarse 打分每条 query 的结构复杂度为 $O(Kd)$，关系分支在分块计算
下为 $O(K_r^2d)$，而不是对全部 256 个候选做二次计算。这是结构复杂度说明。仓库同时提供可配置
候选截断、warm-up、测量批次、batch size 和设备的 forward-only benchmark；它排除数据读取且不计算
检索指标。仓库保留这些脚本作为实现参考；本文不声称 timing、显存、吞吐或部署性能结果。

\section{实验}

\subsection{数据集与候选池}

我们使用 MassSpecGym 论文发布的划分 \cite{bushuiev2024massspecgym}，包含 194,119 条训练
谱图、19,429 条验证谱图和 17,556 条测试谱图，共计 231,104 条谱图；基准论文报告约
29,000 个唯一分子结构。其相似性感知、结构不重叠构造减少了训练和测试分子之间的近邻
结构重叠。

实验评估 MassSpecGym 提供的两种候选设置。质量约束候选池使用前体质量和容差窗口；
分子式约束候选池假设已知分子式，因此应解释为 formula-conditioned retrieval，而不是
完全开放的未知结构注释。冻结候选池至多包含 256 个候选。reranker 是闭集方法，只能
改善所提供候选的顺序，不能恢复候选池之外的分子。

主要评价是 overlap-clean 的 alignment-42 checkpoint 上的 cache-local 矩阵，复用保存的候选
顺序、嵌入和参考二维 InChIKey 身份标签；它不是重新通过官方 loader 编码或重训 alignment。
另有独立的 official-compatible 候选顺序和身份审计，作为协议核对而不是新的主结果。本地
exact-target-SMILES 结果作为敏感性视图。完整协议、身份规则和候选顺序说明见 Appendix~A。

\subsection{实现细节}

对齐模型输出 512 维嵌入。谱图 Transformer 为 4 层、16 个 attention heads、宽度 512；
分子 GINE 为 4 层、隐藏宽度 128、dropout 0.2。当前实验不加载已发布的 SpecEmbedding
checkpoint，因此跳过仅冻结预训练谱图塔的第一训练阶段，并在后续 alignment 阶段使用
MassSpecGym 训练集端到端优化两个编码塔。alignment 使用 AdamW，batch size 128，学习率
$10^{-4}$，weight
decay $10^{-4}$，cosine schedule，gradient-norm clipping 为 1，最多 100 个 epoch，early
stopping patience 为 5。每个 epoch 对每个唯一分子标签采样 10 次，每次随机选择一条关联
谱图。谱图和分子图增强分别以 0.5 概率启用：谱图增强删除归一化强度低于 0.3 的碎片峰中
20\%（向下取整），并进行 $\pm15\%$ 强度扰动；图增强执行 10\% node dropout 和 10\%
directed edge-entry dropout。验证关闭增强，但在固定随机流下保留关联谱图的随机采样。
当前主 reranker 矩阵使用 seed 42、43、44 产生的 alignment checkpoint。每个 checkpoint 内，
reranker seed 42--44 控制初始化、DataLoader 顺序、候选打乱和 dropout；这些 checkpoint 是
内部描述性运行，并不等价于来自已知总体的独立抽样。

reranker 隐藏维度为 512，相对分支维度为 32，pair chunk 为 16，dropout 为 0.1。模型先对
$K=256$ 的完整候选池打分，再对 coarse top-$K_r=40$ 使用相对分支。优化器为 AdamW，学习率
$10^{-4}$，batch size 32，最多训练 30 个 epoch，并根据验证集 MRR 使用 patience 为 5 的
early stopping。pairwise 与谱图依赖损失的权重和 margin 为
$\lambda_{\mathrm{pair}}=0.2$、$\lambda_{\mathrm{spec}}=0.1$、$\gamma=0.1$ 和
$\delta=0.1$。参数近似匹配的 pointwise 对照使用相同绝对分支和优化协议，但不做相对聚合；
历史 candidate-set Transformer 仅作为历史参考。

每项结果工件都列出 alignment checkpoint 和 reranker training seeds。training seed 同时控制
初始化、DataLoader 顺序、候选打乱和 dropout，并不是独立的 alignment 重复。训练只使用正例
自然存在于完整候选池中的 query；无此标签的 query 从 reranker loss 中排除，但保留在覆盖率
统计中。任何 cache 都不插入正例。模型选择不访问测试集。每项正式 reranker 运行最多请求
20,000 条训练 query，并使用三个 reranker seed。三个 alignment checkpoint 和九组
alignment--reranker 配对提供描述性稳健性信息，但不是置信区间、显著性检验或完整端到端
seed 分析。

\subsection{代码、数据与工件可用性}

匿名补充包提供源码检查、锁定的环境文件、自动化测试和一个小型 CPU 合成 smoke test。
其中不包含真实谱图、候选池、训练 checkpoint 或 cache。仓库的复现索引记录正式结果所用的
本地 checkpoint、解析后的配置、日志、哈希和选定的 ranking 输出。真实数据入口及更多复现
边界见 Appendix~D。仓库源码还包含 query-level case exporter，可为指定 cache query 导出候选身份
以及 base/reranker 排名，供后续人工核验；它不执行 fragment attribution，也不生成化学因果解释。

\subsection{对比方法与指标}

受控比较包括余弦基础检索器、参数近似匹配且无 rank 的 pointwise coarse 打分器和本文提出的
相对 reranker；早期 candidate-set Transformer 仅作为历史参考。我们还运行了与 JESTR 类似打分形式的
本地保存嵌入 cosine 对照；该对照是 reimplemented control，不是官方 JESTR 编码器或 checkpoint
复现。由于该对照与基础检索器都在相同的保存归一化嵌入上使用 cosine 分数，二者的本地数字
一致；因此它是协议匹配的打分控制，而不是独立的 JESTR 比较。GLMR 仍为 reported-only，
因为其条件生成流程未在本文中复现。

评价指标为全部 17,556 条测试 query 上的 Recall@1、Recall@5、Recall@10、Recall@20 和 MRR。
主要表使用 alignment seed 42 的 cache-local no-forcing 候选顺序和身份标签、三个 reranker
training seeds（42--44）以及最多 20,000 条训练 query；另一个表报告 alignment seeds 42--44
的相同 pointwise/relative 对照。均值和样本标准差在单个 alignment 内描述 reranker 离散度，
跨 alignment 行是 checkpoint 敏感性描述，不是总体重复。本文不声称官方 JESTR 复现、统计显著性
或端到端部署性能。

\subsection{主要结果}

\begin{table}[t]
  \centering
  \caption{overlap-clean alignment-42 checkpoint 上的主要 cache-local no-forcing 评价。数值为
  百分比；学习式结果为 reranker training seeds 42--44 的均值 $\pm$ 样本标准差。所有方法使用
  同一保存候选池和二维 InChIKey 身份标签；该表不是重新通过 official loader 编码的结果。扩展的
  alignment 敏感性见表~\ref{tab:alignment_sensitivity_cn}。}
  \label{tab:main_results_cn}
  \scriptsize
  \setlength{\tabcolsep}{2pt}
    \begin{tabular}{llrrrrr}
    \toprule
    候选池 & 方法 & R@1 & R@5 & R@10 & R@20 & MRR \\
    \midrule
    \multirow{3}{*}{质量}
      & \base{}（256 候选池） & 47.46 & 63.51 & 71.02 & 77.76 & 55.03 \\
      & Pointwise 对照
      & 58.12$\pm$0.54 & 73.61$\pm$0.54 & 78.63$\pm$0.48 & 82.24$\pm$0.10 & 64.91$\pm$0.47 \\
      & 相对候选模型
      & 57.88$\pm$0.59 & 73.55$\pm$0.58 & 78.72$\pm$0.29 & 82.15$\pm$0.13 & 64.78$\pm$0.52 \\
    \midrule
    \multirow{3}{*}{分子式}
      & \base{}（256 候选池） & 63.10 & 75.79 & 80.87 & 84.93 & 68.86 \\
      & Pointwise 对照
      & 70.76$\pm$0.59 & 80.07$\pm$0.35 & 84.33$\pm$0.08 & 86.86$\pm$0.08 & 75.17$\pm$0.41 \\
      & 相对候选模型
      & 70.88$\pm$0.38 & 80.60$\pm$0.37 & 84.68$\pm$0.26 & 86.85$\pm$0.09 & 75.38$\pm$0.27 \\
    \bottomrule
  \end{tabular}
\end{table}

\noindent 两种学习式模型在两个候选池都超过 alignment-42 的 Base。Pointwise 在质量场景的 R@1
和 MRR 略高于 relative（58.12 对 57.88、0.6491 对 0.6478），而 relative 在分子式场景略高
（70.88 对 70.76、0.7538 对 0.7517）。这些描述性差异不能建立相对模块的独立收益。仓库保留
了评价时候选 cutoff 文件，但它们不是按候选规模重新训练的模型，也不构成额外的主要结果。
seed-42 paired comparison 显示第二阶段相对 Base 有明显改善，但不能据此建立相对模块的独立收益。

这一受控比较形成本文的两部分发现：在所提供候选池内，监督式第二阶段打分确实有用；但当前
观测到的改善并不要求相对交互，而且 relative 并未在匹配的 pointwise 对照上普遍占优。后一个
结果是所提机制的适用边界，而不是需要隐藏的比较失败。

local exact-target-SMILES 敏感性结果见 Appendix~A，不改变上述结论。

\begin{table}[t]
  \centering
  \caption{扩展无 forcing top-40 reranker cache 下的 alignment checkpoint 描述性敏感性。
  学习式结果为 reranker seeds 42--44 的均值 $\pm$ 样本标准差；该矩阵未重新通过 official loader 编码。}
  \label{tab:alignment_sensitivity_cn}
  \scriptsize
  \setlength{\tabcolsep}{2pt}
  \begin{tabular}{@{}llrrr@{}}
    \toprule
    alignment & 候选池 & Base R@1/MRR & Pointwise R@1/MRR & Relative R@1/MRR \\
    \midrule
    42 & 质量 & 47.46/.5503 & 58.12/.6491 & 57.88/.6478 \\
    43 & 质量 & 35.92/.4433 & 58.32/.6398 & 57.61/.6350 \\
    44 & 质量 & 23.91/.3226 & 46.42/.5238 & 47.05/.5284 \\
    \midrule
    42 & 分子式 & 63.10/.6886 & 70.76/.7517 & 70.88/.7538 \\
    43 & 分子式 & 45.08/.5185 & 64.72/.6836 & 64.10/.6795 \\
    44 & 分子式 & 28.73/.3687 & 54.26/.5877 & 54.35/.5911 \\
    \bottomrule
  \end{tabular}
\end{table}

在九组 alignment--候选池--seed 配对中，relative 减 pointwise 的 R@1 平均差在质量和分子式
场景分别为 $-0.11$ 和 $-0.14$ 个百分点，MRR 平均差分别为 $-0.0005$ 和 $+0.0005$；两种
候选协议中 relative 均在 9 组配对中的 5 组取得较高 R@1。这些是描述性配对估计，不是显著性
检验，支持的结论是监督重排序框架比任何独立的 relative 模块优势更稳妥。

\paragraph{机制分析。}
表~\ref{tab:relative_ablations_cn} 报告完整的 30 次 relative 模块消融矩阵：5 种 relative
配置、两种候选协议以及三个 reranker seed，均在 alignment-42 checkpoint 上运行。完整配置为
带谱图条件的反对称模型。改为 directed pair、移除 pair loss、移除谱图 loss 或移除谱图条件后，
两个候选协议的变化方向并不一致；没有一个变体在所有指标上统一超过完整模型。因此该实验
说明的是模块行为边界，而不是组件因果排序。

\begin{table}[t]
  \centering
  \caption{cache-local 无 forcing top-40 协议下的 relative 模块消融。数值依次为 R@1 和 MRR；
  每项为 reranker seeds 42--44 的均值 $\pm$ 样本标准差。}
  \label{tab:relative_ablations_cn}
  \scriptsize
  \setlength{\tabcolsep}{2pt}
  \begin{tabular}{@{}lrr@{}}
    \toprule
    变体 & 质量 R@1/MRR & 分子式 R@1/MRR \\
    \midrule
    完整反对称 & 57.88$\pm$0.59 / .6478$\pm$.0052 & 70.88$\pm$0.38 / .7538$\pm$.0027 \\
    Directed pair & 57.97$\pm$0.49 / .6471$\pm$.0047 & 70.06$\pm$1.29 / .7471$\pm$.0106 \\
    移除谱图条件 & 58.42$\pm$0.81 / .6522$\pm$.0069 & 71.01$\pm$0.40 / .7546$\pm$.0033 \\
    移除 $\mathcal{L}_{\mathrm{pair}}$ & 58.01$\pm$0.38 / .6488$\pm$.0037 & 70.91$\pm$0.36 / .7541$\pm$.0027 \\
    移除 $\mathcal{L}_{\mathrm{spec}}$ & 58.44$\pm$0.54 / .6516$\pm$.0064 & 70.59$\pm$0.79 / .7515$\pm$.0056 \\
    \bottomrule
  \end{tabular}
\end{table}

JESTR 和 GLMR 保留为 reported-only 的相关工作背景；本文不据此推断相对排名。

\paragraph{困难候选分层。}
在 seed-42 本地 query 分析中，选定的有标签子集包含质量 14,726 条和分子式 15,670 条 query。
relative 模型在基础排名为 2--5/6--20/21--40 时相对 Base 的 R@1 增益，质量场景分别为
+54.05/+38.13/+29.98 个百分点，分子式场景分别为 +55.41/+34.18/+19.10 个百分点。基础
排名为 1 的 query 则下降到质量 87.41\%、分子式 95.12\%，而 Base 分层为 100\%。Pointwise
对照也呈现相同的未解决 query 模式；这些分层是描述性分析，不建立相对模块的因果效应。

早期含 rank 的 Transformer 结果和协议说明移至 Appendix~B，不作为当前相对方法的性能证据。

\section{讨论}

扩展 pilot 将候选覆盖与列表内排序分开。三个 alignment checkpoint 的每个 checkpoint 内，
两种模型都超过对应 Base；pointwise 在主要 R@1/MRR 汇总略强，relative 在不同 cutoff 并非
始终更好。该结果支持监督第二阶段重排序作为框架组件，但不能隔离相对模块的独立收益。

30 次消融显示，directed 与 antisymmetric pair、谱图条件以及两项辅助损失对两个候选池的
影响方向不同。这些结果是描述性分析，用于说明方法行为，不构成组件因果收益。

relative 分支的结构计算成本也高于 pointwise：完整候选池的绝对分支为 $O(Kd)$，coarse 子集
上的 pair 计算为 $O(K_r^2d)$。仓库保留固定硬件的 forward-only 测量作为实现参考；本文不将其
解释为端到端 latency、吞吐或部署性能结论。

主要矩阵复用 cache-local 候选顺序和保存的 alignment 嵌入；另有独立的 official-compatible 审计
使用发布的候选 JSON 和相同的保存嵌入，但没有通过新建的 official loader 重新编码谱图/分子，
也没有重训第一阶段。因此它不能证明另行重建的 loader、不同候选来源或外部方法的端到端等价性。

JESTR/GLMR 仍为 reported-only，未独立复现。

\section{局限性与负责任评估}

扩展 reranker 矩阵最多请求 20,000 条训练 query，评价三个 alignment checkpoint，并在每个
主要配置上使用三个 reranker seed。这提供描述性 checkpoint 敏感性，而不是置信区间、显著性
检验或完整端到端 seed 分析。评价复用保存嵌入和所提供候选池，不声称重新通过官方 loader
编码、外部 baseline 复现或候选感知 alignment 重训。这些限制也界定了本文的两部分发现：
框架级重排序观察在所评价 checkpoint 内得到支持，而 relative 模块的普适性仍未确定。

本文方法是闭集方法，在分子式协议下还假设分子式已知。它不能恢复候选池之外的目标，也不能
自动迁移到其它候选构造和采集条件。谱图塔使用选定峰和前体 $m/z$，但不显式接收加合物、
仪器或碰撞能量元数据。其 pairwise 结构约束也不提供概率校准或全局一致排序保证。更多协议
和复现边界见 Appendix~A 和 Appendix~D。

\section{结论}

本文提出 \method{}：一种无 rank、无 forcing 的第二阶段排序器，将完整候选池 coarse 打分
与谱图条件的相对候选偏好结合。在扩展的 cache-local 矩阵中，三个 alignment checkpoint 内的
relative 和参数近似匹配的 pointwise 都改善质量和分子式场景的对应 Base；主要表使用
overlap-clean 的 alignment-42 保存嵌入候选池，另有独立审计核对发布的候选顺序和二维身份规则。
Pointwise 在主要 R@1/MRR 汇总略强，30 次消融也显示 relative 组件没有统一的独立优势。因此，
当前证据支持一个边界明确的监督式重排序框架，并明确了 relative 交互的边界：它是具有结构依据
的候选关系机制，但当前证据尚未建立其相对 pointwise 打分的独立且普遍稳健收益。本文方法仍
限定于固定的保存嵌入候选协议。

\appendix

\section{评价协议与候选身份}

主要多 seed 矩阵是 cache-local 评价：在 overlap-clean 的 alignment-42 checkpoint 上复用保存的
候选顺序、嵌入和参考二维 InChIKey 身份标签。另有独立的 official-compatible 审计使用
MassSpecGym 1.3.1 retrieval 候选顺序和相同的保存嵌入；它不重新通过 official loader 编码，也不
重训 alignment。两种范围保持区分，因此主表不被表述为一次全新的官方评价。

受跟踪的完整池身份检查在两个候选协议的 17,556 条测试 query 上均找到一个身份等价正例，
没有多正例 query 或候选身份碰撞。因此，在这些保存候选列表上，exact-target 标签与参考身份
标签等价。本地 exact-target-SMILES 结果作为敏感性视图保留。质量候选使用前体质量容差；
分子式候选假设公式已知。两种设置都是闭集方法，目标不在评价截断内时 reranker 不能恢复它。

谱图记录签名分析还记录了 6 条验证和 3 条测试记录与训练记录匹配、但标签分子不同的情况。
overlap-clean 配置在验证阶段排除受影响记录；用于敏感性分析的三个 alignment checkpoint
只是内部描述性 estimates，不是总体稳定性声明。评价时的候选截断敏感性使用同一 256-candidate
cache，并非按候选规模重新训练；保留的 cutoff 文件是辅助工件，不构成额外的主表结果。

\section{实现细节与历史对照}

alignment 模型输出 512 维归一化嵌入。谱图塔为 4 层、16 heads、宽度 512 的 Transformer；
分子塔为 4 层、隐藏宽度 128 的 GINE。reranker 隐藏宽度为 512、关系宽度为 32、pair chunk
为 16、dropout 为 0.1，使用 $K=256$ 和 $K_r=40$。两阶段均使用 AdamW；reranker 使用
$10^{-4}$ 学习率、batch size 32、验证集 MRR early stopping，以及正文所列损失权重和 margin。

reranker seed 控制初始化、DataLoader 顺序、候选打乱和 dropout。42--44 的均值和样本标准差
描述 reranker 层级离散度，不是独立 alignment 重复。当前训练中自然没有正例的 query 保留在
覆盖率统计中，但不进入 loss；当前方法 cache 不插入正例。

当前训练入口会在新 reranker checkpoint 中写入轻量的 \texttt{data\_summary}，记录实际有标签训练
query 数、请求的最大 query 数、\texttt{train\_k}、model type、pair mode、seed，以及可用的 cache
协议字段（如 split、candidate type、预检索 cutoff 和 positive-forcing 状态）。这些元数据用于区分
pilot 工件和其它训练范围，不进入 forward，也不能替代完整训练集、alignment 层级或不确定性实验。

历史 Candidate-Set Transformer 使用 rank embedding 和通用候选 self-attention，其训练期
positive forcing 和 rank-aware 配置不属于当前 relative/pointwise 正式协议。它只用于连接早期
设计，且参数量和训练协议不同，不能作为容量匹配的主要对照。

\section{补充机制分析}

表~\ref{tab:relative_ablations_cn} 的 30 次消融覆盖完整反对称模型、directed pair、移除谱图
条件、移除 pair loss 和移除谱图 loss，横跨两个候选协议与 reranker seeds 42--44。没有任何
消融在所有指标上统一最优，因此表中结果是关于模块行为和敏感性的证据，而不是组件因果估计。

困难候选分层显示，relative 相对 Base 的主要增益集中在基础排名为 2--5 和 6--20 的 query，
而基础排名为 1 或高相似度 query 可能退化。这一模式支持候选关系假设，但不表示关系模块
对所有 query 都有益。

上述困难候选分层也可视为固定候选池剩余空间的检查：增益主要出现在基础排名 2--5 和
6--20，而排名 1 query 可能退化。评价时的 cutoff 结果和固定硬件前向成本只作为实现参考，
不是端到端部署基准。

\section{复现与负责任评价说明}

正式矩阵最多请求 194,119 条可用训练 query 中的 20,000 条，使用 alignment checkpoints 42--44，
并以 seeds 42--44 重复 reranker 配置。复现索引记录 source commit、解析后的配置、cache 元数据、
checkpoint 和日志哈希，以及选定的 ranking 输出。匿名补充包包含源码检查、锁定环境、测试和合成
CPU smoke test，但不包含真实谱图、候选池、cache 或训练 checkpoint，因此 smoke test 不能复现科学
指标；本地工件边界在 Appendix~D 中明确说明。

本地保存嵌入 cosine 对照是 JESTR-style scoring 的 reimplemented control，不是官方 JESTR
编码器或 checkpoint 复现；GLMR 仍为 reported-only。本文不声称 SOTA、统计显著性、完整
alignment 稳定性、跨数据集泛化或端到端部署 latency；这些内容界定方法结果的适用范围，而
不是额外实验结论。

\bibliographystyle{plain}
\bibliography{references}

\end{document}
